% =============================================================
% Abstrak, based on Lampiran 10 of the Pedoman.
% One paragraph only, 200-300 words: latar belakang, tujuan,
% metodologi, dan hasil. No tables/figures/direct quotations.
% =============================================================
\thispagestyle{empty}
\begin{center}
{\bfseries ABSTRAK\par}
\vspace{10mm}
\end{center}

\noindent Perkembangan volume data yang pesat pada berbagai sektor menuntut pendekatan analisis yang mampu mengekstraksi pola dan pengetahuan secara otomatis dan akurat. Penelitian ini dilatarbelakangi oleh permasalahan tersebut, khususnya pada domain \textit{[nama domain aplikasi, misalnya: deteksi anomali/klasifikasi teks/prediksi deret waktu]}, di mana pendekatan konvensional belum mampu memberikan hasil yang optimal dari segi akurasi maupun efisiensi komputasi. Penelitian ini bertujuan untuk merancang dan mengevaluasi sebuah pendekatan berbasis \textit{[nama metode/algoritma, misalnya: pembelajaran mendalam, model probabilistik Bayesian, atau algoritma metaheuristik]} yang dapat meningkatkan kinerja pada permasalahan tersebut dibandingkan dengan metode \textit{baseline}. Metodologi penelitian terdiri atas tahap pengumpulan dan praproses data, perancangan arsitektur model, pelatihan dan penyetelan hiperparameter, serta pengujian menggunakan skema validasi silang pada himpunan data \textit{[nama/asal dataset]}. Kinerja model dievaluasi menggunakan metrik \textit{[nama metrik, misalnya: akurasi, F1-score, RMSE]} dan dibandingkan dengan beberapa metode pembanding dari literatur terkini. Hasil penelitian menunjukkan bahwa metode yang diusulkan memberikan peningkatan kinerja sebesar \textit{[nilai/persentase]} dibandingkan dengan metode pembanding, serta menunjukkan konsistensi hasil pada beberapa skenario pengujian. Temuan ini mengindikasikan bahwa pendekatan yang diusulkan berpotensi diterapkan pada permasalahan sejenis dalam skala yang lebih luas, sekaligus membuka peluang penelitian lanjutan terkait optimalisasi arsitektur model dan efisiensi komputasi.

\vspace{6mm}
\noindent\textbf{Kata Kunci:} kata kunci satu, kata kunci dua, kata kunci tiga, kata kunci empat.
\clearpage
